\chapter{Asymmetric Volatility}
\label{ch:asymmetry-jumps}

Layer~1 decomposes realized variance into signed components and jump variation.
Six features, all computed from tick-level 5-minute returns.

%% ──────────────────────────────────────────────
\section{Features}
\label{sec:asymmetry-features}
%% ──────────────────────────────────────────────

\begin{table}[htbp]
\centering
\caption{Layer~1 feature set: asymmetric volatility and jump decomposition.}
\label{tab:asymmetry-features}
\small
\begin{tabular}{@{}p{2.8cm}p{4.2cm}p{4.8cm}l@{}}
\toprule
\textbf{Feature} & \textbf{What It Is} & \textbf{What It Does} & \textbf{Horizon} \\
\midrule
$\mathrm{RS}_t^{-}$ (daily) &
  $\displaystyle\sum_{i=1}^{M} r_{t,i}^{2}\,\mathbf{1}(r_{t,i} < 0)$ &
  Carries ${\sim}2\times$ predictive weight of $\mathrm{RS}^{+}$; dominant asymmetric signal \citep{PattonSheppard2015} &
  1d \\[10pt]
$\mathrm{RS}_t^{+}$ (daily) &
  $\displaystyle\sum_{i=1}^{M} r_{t,i}^{2}\,\mathbf{1}(r_{t,i} \geq 0)$ &
  Weaker predictor; provides the contrast that identifies asymmetry &
  1d \\[10pt]
$\mathrm{RS}_t^{-(w)}$ (weekly) &
  $\displaystyle\frac{1}{5}\sum_{j=0}^{4} \mathrm{RS}_{t-j}^{-}$ &
  Persistent downside memory; smooths daily noise in the negative semivariance &
  1d--5d \\[10pt]
$J_t^{-}$ (signed neg.\ jump) &
  Large negative moves beyond threshold &
  1--3\% QLIKE gain beyond unsigned jumps \citep{ABD2007} &
  1d--5d \\[10pt]
$C_t$ (continuous variation) &
  $\max(\BPV_t,\; 0)$ &
  Highly persistent (ACF ${\sim}0.6$--$0.7$); workhorse of the decomposition \citep{BNS2004} &
  All \\[10pt]
$J_t$ (jump variation) &
  $\max(\RV_t - \BPV_t,\; 0)$ &
  Nearly unpredictable (ACF ${\sim}0.0$--$0.1$); signals regime breaks &
  Event \\
\bottomrule
\end{tabular}
\end{table}


%% ──────────────────────────────────────────────
\section{The Leverage Effect}
\label{sec:leverage-effect}
%% ──────────────────────────────────────────────

Negative returns increase future volatility more than positive returns of equal magnitude.
The effect is strongest for equity indices and E-mini futures, where $\mathrm{RS}^{-}$ carries roughly twice the predictive weight of $\mathrm{RS}^{+}$ in HAR-type regressions \citep{PattonSheppard2015}.
The SHAR model exploits this directly by replacing $\RV$ with separate $\mathrm{RS}^{+}$ and $\mathrm{RS}^{-}$ regressors at each horizon.


%% ──────────────────────────────────────────────
\section{Continuous vs.\ Jump Variation}
\label{sec:continuous-vs-jump}
%% ──────────────────────────────────────────────

Continuous variation $C_t$ is highly persistent and drives the bulk of out-of-sample forecast accuracy; it behaves like a smoothed version of $\RV$ with jump contamination removed.
Jump variation $J_t$ is nearly unpredictable day-to-day, but large positive values flag regime transitions (macro announcements, flash crashes) where the conditional distribution shifts.
In practice, $C_t$ is the forecasting workhorse and $J_t$ is the alarm: $C$ enters the model as a persistent regressor at all horizons, while $J$ triggers recalibration or regime-switch logic.
The HAR-CJ model of \citet{ABD2007} formalizes this split and shows that separating $C$ from $J$ yields 1--3\% QLIKE improvement over models that use total $\RV$ alone.


%% ──────────────────────────────────────────────
\section{Cumulative Performance}
\label{sec:cumulative-l1}
%% ──────────────────────────────────────────────

Layers~0+1 together (11 features: 5 from HAR/HARQ, 6 from this chapter) achieve approximately 70\% of the total attainable QLIKE improvement, per literature benchmarks.
All semivariances and jump components are computed from tick-level data at 5-minute frequency.
Jump significance is assessed via the Lee-Mykland test, which flags individual intraday returns exceeding a time-varying threshold calibrated to local bipower variation.


%% ──────────────────────────────────────────────
%% Boxes
%% ──────────────────────────────────────────────

\begin{keyidea}[Signed Semivariances: The Cheapest Upgrade]
Replacing $\RV$ with $\mathrm{RS}^{+}$ and $\mathrm{RS}^{-}$ in a HAR regression (the SHAR model) costs zero additional data and yields 3--8\% QLIKE improvement \citep{PattonSheppard2015}.
This is the single cheapest feature upgrade after the RQ interaction from Layer~0.
\end{keyidea}

\begin{warning}[Jumps Are Not Predictors]
$J_t$ has near-zero autocorrelation.
Do not include lagged jump variation as a standard regressor expecting it to forecast future $\RV$.
Its value is as a conditioning variable: when $J_t$ is large, shrink the daily coefficient or trigger a regime indicator.
Treating $J_t$ as a regular feature wastes a degree of freedom and can degrade forecast accuracy.
\end{warning}
